\subsection{SAT Solver (Using SAT Library)}
\subsubsection{Idea}
The SAT approach transforms the Futoshiki puzzle into a propositional satisfiability problem in CNF (Conjunctive Normal Form). Instead of searching directly over puzzle states, we encode all puzzle rules as logical clauses and let a SAT engine decide whether there exists a satisfying assignment.

In this project, CNF clauses are generated by the \texttt{FOLKB} module, then solved by PySAT (Glucose3). If SAT, the returned model is decoded back to an \(N \times N\) grid; if UNSAT, the puzzle has no valid solution under given constraints.

\subsubsection{Constraint Encoding}
For each triple \((i,j,v)\), define boolean variable \(X_{i,j,v}\):
\[
X_{i,j,v} = \text{true} \iff \text{cell }(i,j)\text{ has value }v.
\]

Variable ID mapping used in code:
\[
\text{ID}(i,j,v)= i\cdot N^2 + j\cdot N + v,
\]
with \(i,j\in[0,N-1]\), \(v\in[1,N]\), so IDs are in \([1,N^3]\).

Main CNF groups:
\begin{itemize}
	\item \textbf{Cell at least one value:} each cell must take some \(v\in\{1,\dots,N\}\).
	\item \textbf{Cell at most one value:} no cell can take two different values simultaneously.
	\item \textbf{Row uniqueness:} two different columns in same row cannot share the same value.
	\item \textbf{Column uniqueness:} two different rows in same column cannot share the same value.
	\item \textbf{Given clues:} pre-filled cells are encoded as unit clauses.
	\item \textbf{Inequality constraints:} for each adjacent pair with sign \(<\) or \(>\), all violating value-pairs are forbidden by binary clauses.
\end{itemize}

This encoding guarantees that any satisfying model corresponds to a valid Futoshiki solution.

\subsubsection{Pseudocode}
\begin{algorithm}[H]
\caption{SAT-based Futoshiki Solver}
\begin{algorithmic}[1]
\State Build CNF clauses from puzzle constraints using FOLKB
\State Create SAT solver instance (Glucose3)
\State Add all CNF clauses to solver
\If{solver returns UNSAT}
	\State \Return \textbf{failure}
\EndIf
\State Get SAT model (set of literals)
\State Initialize empty grid \(N \times N\)
\For{each positive literal \(\ell\) in model}
	\If{\(1 \le \ell \le N^3\)}
		\State Decode \(\ell \rightarrow (i,j,v)\)
		\State Set grid\((i,j) \leftarrow v\)
	\EndIf
\EndFor
\If{any cell remains unassigned}
	\State \Return \textbf{failure}
\Else
	\State \Return solved grid
\EndIf
\end{algorithmic}
\end{algorithm}

\subsubsection{Illustration}
For a \(4 \times 4\) puzzle:
\begin{enumerate}
	\item The encoder creates variables \(X_{i,j,v}\) for all \(4^3=64\) combinations.
	\item Clauses are generated for Latin constraints (row/column uniqueness), given clues, and inequality symbols.
	\item Glucose3 finds a satisfying model, for example including literals that decode to assignments such as \((0,0)=2\), \((0,1)=4\), etc.
	\item After decoding all positive in-range literals, the grid is fully reconstructed and validated as complete.
\end{enumerate}

Compared with incremental state expansion, the SAT method delegates global consistency reasoning to a highly optimized conflict-driven SAT engine.

\subsubsection{Algorithm Analysis}
\textbf{Encoding cost:}
\begin{itemize}
	\item Number of variables is \(N^3\).
	\item Number of clauses is polynomial in \(N\), dominated by pairwise exclusivity and inequality forbidding constraints.
\end{itemize}

\textbf{Solving cost:}
\begin{itemize}
	\item SAT is NP-complete in general, so worst-case complexity is exponential.
	\item In practice, modern CDCL solvers (such as Glucose3) are often very effective on structured constraints like Futoshiki.
\end{itemize}

\textbf{Space and implementation trade-offs:}
\begin{itemize}
	\item Requires storing full CNF and SAT internal structures.
	\item Provides a clean and extensible formulation: adding new logical constraints usually means adding new clause templates.
\end{itemize}

\textbf{Practical behavior in this project:}
\begin{itemize}
	\item SAT solver gives deterministic feasibility checking (SAT/UNSAT).
	\item Performance is strong on many instances, especially when constraints are dense enough to guide conflict learning.
\end{itemize}
